\section{Conclusion and Future Work}\label{sec:conclusion}

This work demonstrates that physics-gated learning is a practical path to resilient orbital correction in the Solar Cycle 25 era. By hard-constraining the residual model with $\Gamma(t)=\tanh(\lambda t)$, the G-PINN removes Phantom Drift at epoch while preserving the efficiency and interoperability of SGP4. Across storm regimes, the framework improves mean prediction accuracy, compresses extreme error tails (28 km to 21 km), and supports tighter operational search volumes.

A second outcome is diagnostic: the persistent $+5.5$ m/s radial correction suggests a systematic bias in baseline drag handling for selected public-element tracks. This indicates that physics-informed residual models can serve not only as predictors but also as quality-control layers for legacy propagation chains.

Operationally, integrating G-PINN corrections with Monte Carlo covariance propagation, Mahalanobis risk classes, and WebGL inspection tools enables end-to-end decision support for conjunction response during severe geomagnetic activity.

Future work focuses on: (i) constellation-scale graph learning to share drag representations across orbital shells, (ii) direct ingestion of high-cadence onboard telemetry to reduce TLE latency effects, and (iii) coupling with predictive space-weather models for proactive rather than reactive drag correction.